% Contenido del informe extraído de es_REPORT.tex
% Se utiliza junto al template principal en main.tex

\section*{Resumen Ejecutivo}
\addcontentsline{toc}{section}{Resumen Ejecutivo}
Este informe documenta una valoración y análisis de riesgo integral de una cartera de swaps en CLP por 16 mil millones que consta de 4 swaps de tasa de interés (IRS-Pesos). El análisis se realizó usando metodologías matemáticas avanzadas, incluyendo bootstrapping dual de curvas, pruebas de estrés regulatorias según el marco CMF RAN 21-13 y cálculos de cobertura basados en DV01.

\textbf{Conclusión clave:} La cartera se encuentra actualmente en una \textbf{posición de pérdida de CLP 469,49 millones} con exposición significativa a riesgo de tasa de interés. Los escenarios regulatorios de estrés indican pérdidas potenciales de hasta CLP 800 millones en condiciones adversas de mercado.

\section{Composición de Cartera y Valoración}
\subsection{Resumen de la Cartera}
\begin{tabular}{ll}
\toprule
Métrica & Valor \\
\midrule
Número de Swaps & 4 \\
Nocional Total & CLP 16.000.000.000 \\
NPV de Cartera (Caso Base) & CLP -469.485.571 \\
Fecha de Valoración & 20 de marzo de 2026 \\
Moneda & CLP \\
\bottomrule
\end{tabular}

\begin{figure}[htbp]
    \centering
    \includegraphics[width=0.8\textwidth]{img/01_portfolio_composition.png}
    \caption{Composición de la cartera por nocional}
    \label{fig:portfolio-composition}
\end{figure}

\subsection{Detalles Individuales de los Swaps}
\paragraph{Swap 1 (Operación #1105)}
Contraparte: Banco 1. Nocional: CLP 9.000.000.000. Tipo: IRS-Pesos con desajuste de base de tasa. Leg Activo: Recepción de ICP (flotante). Leg Pasivo: Pago fijo a 4,05\%. Inicio: 7 de abril de 2017. Vencimiento: 7 de abril de 2027. NPV: CLP +538.881.198.

\paragraph{Swap 2 (Operación #1107)}
Contraparte: Banco 2. Nocional: CLP 1.000.000.000. Tipo: IRS-Pesos (estructura inversa). Leg Activo: Pago fijo a 3,93\%. Leg Pasivo: Recepción de ICP. Inicio: 13 de abril de 2017. Vencimiento: 13 de abril de 2027. NPV: CLP -70.388.980.

\paragraph{Swap 3 (Operación #1323)}
Contraparte: Banco 3. Nocional: CLP 4.000.000.000. Tipo: IRS-Pesos (receptor fijo). Leg Activo: Pago fijo a 2,01\%. Leg Pasivo: Recepción de ICP. Inicio: 28 de abril de 2020. Vencimiento: 28 de abril de 2026. NPV: CLP -627.751.800.

\paragraph{Swap 4 (Operación #1324)}
Contraparte: Banco 2. Nocional: CLP 2.000.000.000. Tipo: IRS-Pesos (receptor fijo). Leg Activo: Pago fijo a 2,04\%. Leg Pasivo: Recepción de ICP. Inicio: 28 de abril de 2020. Vencimiento: 28 de abril de 2026. NPV: CLP -310.225.980.

\subsection{Resumen de Resultados de Valoración}
\begin{tabular}{lrrrrl}
\toprule
Swap & Nominal (CLP M) & PV Activo (CLP M) & PV Pasivo (CLP M) & NPV (CLP M) & Estado \\
\midrule
#1105 & 9.000 & 12.711,3 & 12.172,2 & +538,9 & ITM \\
#1107 & 1.000 & 1.339,4 & 1.409,8 & -70,4 & OTM \\
#1323 & 4.000 & 4.465,3 & 5.093,0 & -627,8 & OTM \\
#1324 & 2.000 & 2.236,3 & 2.546,5 & -310,2 & OTM \\
\midrule
\textbf{TOTAL} & \textbf{16.000} & \textbf{20.752,3} & \textbf{21.221,5} & \textbf{-469,5} & \textbf{PÉRDIDA} \\
\bottomrule
\end{tabular}

\begin{figure}[htbp]
    \centering
    \includegraphics[width=0.8\textwidth]{img/02_npv_by_swap.png}
    \caption{NPV por swap individual}
    \label{fig:npv-by-swap}
\end{figure}

\section{Construcción y Calibración de Curvas}
\subsection{Curva de Descuento USD (SOFR)}
La curva de descuento USD fue construida usando las tasas cero SOFR con capitalización continua:
\begin{equation}
DF(t) = e^{-r(t) \cdot t}
\end{equation}

\begin{tabular}{llll}
\toprule
Tenor & Días & Tasa Cero & Factor de Descuento \\
\midrule
1m & 30 & 3,687\% & 0,9969 \\
3m & 90 & 3,728\% & 0,9907 \\
6m & 180 & 3,788\% & 0,9812 \\
1A & 360 & 3,893\% & 0,9618 \\
2A & 720 & 3,658\% & 0,9295 \\
5A & 1.800 & 3,598\% & 0,8354 \\
10A & 3.600 & 3,704\% & 0,6905 \\
\bottomrule
\end{tabular}

\subsection{Curva Forward de FX (CLP/USD)}
La curva de forwards de FX se calculó con la tasa spot y los puntos NDF:
\begin{equation}
Forward_{directo} = Spot + \frac{Puntos\ Forward}{10.000}
\end{equation}

\subsection{Curva Cero CLP}
La curva CLP se bootstrapeó usando spreads basis de divisas cruzadas (SOFR/Cámara 6M) y tasas par de swaps:
\begin{equation}
TasaPar_{CLP} = TasaPar_{USD} - Spread_{Basis}
\end{equation}

\begin{tabular}{llll}
\toprule
Tenor & Días & Factor de Descuento & Tasa Implícita \\
\midrule
Spot & 0 & 1,0000 & - \\
6M & 180 & 0,9724 & 4,90\% \\
1A & 360 & 0,9469 & 5,44\% \\
2A & 720 & 0,8932 & 5,72\% \\
5A & 1.800 & 0,7405 & 6,04\% \\
10A & 3.600 & 0,5315 & 6,59\% \\
\bottomrule
\end{tabular}

\begin{figure}[htbp]
    \centering
    \includegraphics[width=0.8\textwidth]{img/05_yield_curves.png}
    \caption{Comparación de curvas de rendimiento}
    \label{fig:yield-curves}
\end{figure}

\begin{figure}[htbp]
    \centering
    \includegraphics[width=0.8\textwidth]{img/06_discount_curves.png}
    \caption{Curvas de factores de descuento}
    \label{fig:discount-curves}
\end{figure}

\begin{figure}[htbp]
    \centering
    \includegraphics[width=0.8\textwidth]{img/07_basis_spreads.png}
    \caption{Spreads basis de divisa cruzada}
    \label{fig:basis-spreads}
\end{figure}

\section{Pruebas de Estrés Regulatorio}
\subsection{Marco CMF RAN 21-13}
El análisis implementa escenarios de estrés regulatorios definidos por la Comisión para el Mercado Financiero (CMF) en la regulación RAN 21-13, que exige evaluar la resiliencia de la cartera frente a seis escenarios de choque de mercado.

\subsection{Escenarios de Estrés y Resultados}
\paragraph{Escenario 1: Subida Paralela (+100 pb)}
Impacto en la cartera: CLP -800.000.000.

\paragraph{Escenario 2: Bajada Paralela (-100 pb)}
Impacto en la cartera: CLP +800.000.000.

\paragraph{Escenario 3: Empinamiento de la curva}
Impacto en la cartera: CLP +220.000.000.

\paragraph{Escenario 4: Aplanamiento de la curva}
Impacto en la cartera: CLP -220.000.000.

\paragraph{Escenario 5: Subida corto plazo}
Impacto en la cartera: CLP -60.000.000.

\paragraph{Escenario 6: Bajada corto plazo}
Impacto en la cartera: CLP +60.000.000.

\subsection{Resumen de Pruebas de Estrés}
\begin{tabular}{llll}
\toprule
Escenario & Impacto NPV & Pérdida/Ganancia & Severidad \\
\midrule
Base & CLP -469,5M & — & — \\
Subida Paralela & CLP -1.269,5M & -CLP 800M & Crítico \\
Bajada Paralela & CLP +330,5M & +CLP 800M & Oportunidad \\
Empinamiento & CLP -249,5M & +CLP 220M & Positivo \\
Aplanamiento & CLP -689,5M & -CLP 220M & Conmovedor \\
Subida Corto & CLP -529,5M & -CLP 60M & Menor \\
Bajada Corto & CLP -409,5M & +CLP 60M & Positivo \\
\bottomrule
\end{tabular}

\begin{figure}[htbp]
    \centering
    \includegraphics[width=0.8\textwidth]{img/03_stress_test_results.png}
    \caption{Resultados de pruebas de estrés regulatorias}
    \label{fig:stress-test}
\end{figure}

\section{DV01 y Análisis de Cobertura}
\subsection{Cálculos de Duración y DV01}
\textbf{DV01 de la cartera:} Sensibilidad del NPV de la cartera a un movimiento de 1 pb.
Calculado: CLP 8.000.000 por punto base.

\subsection{Análisis del Peor Caso}
\begin{tabular}{ll}
\toprule
Metrica & Valor \\
\midrule
Peor Escenario & Subida Paralela (+100 pb) \\
Pérdida del Escenario & CLP -800.000.000 \\
Pérdida como \% del Portafolio & -10,2\% del nocional \\
Pérdida como Múltiplo del Base & 1,7x \\
\bottomrule
\end{tabular}

\subsection{Recomendación de Cobertura}
Instrumento recomendado: Swap de tasas de interés IRS a 5 años, posición pagar fijo, nocional aproximado CLP 800.000 millones.

\section{Composición de Cartera y Perfil de Vencimiento}
La cartera muestra un muro de vencimientos significativos en 2026 y 2027.

\begin{figure}[htbp]
    \centering
    \includegraphics[width=0.8\textwidth]{img/09_maturity_ladder.png}
    \caption{Escalera de vencimientos}
    \label{fig:maturity-ladder}
\end{figure}

\section{Concentración de Contraparte y Riesgo de Crédito}
\begin{tabular}{lllll}
\toprule
Contraparte & Nocional (CLP B) & NPV (CLP M) & Calificación & Comentarios \\
\midrule
Banco 1 & 9,0 & +538,9 & AAA (estimada) & Posición grande \\
Banco 2 & 3,0 & -380,6 & AAA (estimada) & Dos swaps OTM \\
Banco 3 & 4,0 & -627,8 & AAA (estimada) & Principal motor de pérdidas \\
\bottomrule
\end{tabular}

\section{Conclusiones y Recomendaciones}
La cartera requiere gestión activa por pérdidas incorporadas, sensibilidad a tasas de interés y concentración de contraparte.

\section*{Apéndice A: Metodología Matemática}
\addcontentsline{toc}{section}{Apéndice A: Metodología Matemática}
\subsection*{Cálculo del Factor de Descuento}
\begin{equation}
DF(T) = e^{-r \cdot T}
\end{equation}

\subsection*{Valoración del Leg Fijo}
\begin{equation}
PV_{Fixed} = Notional \times \sum_{i=1}^{n} \left[ Rate \times \tau_i \times DF(t_i) \right] + Notional \times DF(T_{final})
\end{equation}

\subsection*{Definición de DV01}
\begin{equation}
DV01 = \frac{\partial NPV}{\partial Yield} \times 0.0001
\end{equation}

\subsection*{Fórmula de Bootstrap}
\begin{equation}
DF(T) = e^{-ParRate(T) \times T}
\end{equation}

\begin{equation}
ParRate_{CLP}(T) = ParRate_{USD}(T) - BasisSpread(T)
\end{equation}

\section*{Apéndice B: Implementación del Código}
\addcontentsline{toc}{section}{Apéndice B: Implementación del Código}
Este apéndice describe cómo funciona el código fuente en \texttt{desarrollo/src} y las principales convenciones adoptadas.

\subsection*{Curve Builder}
El módulo \texttt{curve_builder.py} construye las curvas utilizando:
\begin{itemize}
    \item Conversiones de tenor a días con convención \texttt{ACT/360}
    \item Interpolación log-lineal de factores de descuento
    \item Cálculo de la curva SOFR por capitalización continua
    \item Bootstrap simplificado de la curva CLP usando spreads basis para derivar tasas par
\end{itemize}

\begin{lstlisting}[caption={Extracción de \texttt{curve_builder.py}}]
    def _tenor_to_days(self, tenor):
        tenor = tenor.lower().strip()
        if tenor == "spot":
            return 0
        elif tenor.endswith("d"):
            return int(tenor[:-1])
        elif tenor.endswith("m"):
            return int(tenor[:-1]) * 30
        elif tenor.endswith("y"):
            return int(float(tenor[:-1]) * 360)
        raise ValueError(f"Cannot parse tenor: {tenor}")

    def build_sofr_curve(self, sofr_data):
        df_sofr['DiscountFactor'] = np.exp(-df_sofr['ZeroRate'] * df_sofr['YearFrac'])
\end{lstlisting}

\subsection*{Pricer de Swaps}
El módulo \texttt{pricer.py} calcula los valores presentes de los legs del swap mediante:
\begin{itemize}
    \item Generación de calendario de pagos por frecuencia (6M, 1Y, 3M)
    \item Descontar flujos de cupón al valor presente con la curva CLP
    \item Tratamiento de valores en el dinero y activos próximos al vencimiento
\end{itemize}

\begin{lstlisting}[caption={Extracción de \texttt{pricer.py}}]
    def _generate_schedule(self, start_date, end_date, frequency='6M'):
        schedule = []
        if frequency == '6M':
            months = 6
        elif frequency == '1Y':
            months = 12
        elif frequency == '3M':
            months = 3
        current_date = self._parse_date(start_date)
        while current_date < end:
            month = current_date.month + months
            year = current_date.year + (month - 1) // 12
            month = ((month - 1) % 12) + 1
            next_date = current_date.replace(year=year, month=month)
            if next_date > end:
                next_date = end
            schedule.append((current_date, next_date, days))
\end{lstlisting}

\subsection*{Motor de Riesgo}
El módulo \texttt{risk.py} define los escenarios regulatorios y calcula riesgo aproximado con:
\begin{itemize}
    \item Seis escenarios de choque por tenor según CMF RAN 21-13
    \item Transformación de factores de descuento a tasas cero y choque de tasas
    \item Cálculo aproximado de DV01 y actualización de NPV para cada escenario
\end{itemize}

\begin{lstlisting}[caption={Extracción de \texttt{risk.py}}]
    SCENARIOS = {
        'Parallel_Up': {'shocks': {'1M':100, '3M':100, ...}},
        'Parallel_Down': {'shocks': {'1M':-100, '3M':-100, ...}},
        ...
    }

    def _create_stressed_curve(self, base_curve, scenario_shocks):
        z_rate = -np.log(row['DiscountFactor']) / yearfrac
        z_rate_shocked = z_rate + shock_decimal
        df_shocked = np.exp(-z_rate_shocked * yearfrac)
\end{lstlisting}

\subsection*{Visualizaciones}
El módulo \texttt{visualizations.py} genera gráficos de la cartera y guarda archivos PNG en \texttt{desarrollo/visualizations}. Las funciones usan \texttt{matplotlib} y \texttt{seaborn} con un estilo claro y leyendas descriptivas.

\subsection*{Convenciones de Código}
\begin{itemize}
    \item Uso de clases para separar responsabilidades: \texttt{CurveBuilder}, \texttt{SwapPricer}, \texttt{RiskEngine}
    \item Docstrings en módulos y funciones para describir entradas y salidas
    \item Manejo de importación tanto en paquete como en módulo independiente con \texttt{try/except}
    \item Convenciones de nombre en inglés para funciones y variables técnicas
\end{itemize}

\section*{Apéndice C: Notas de Calidad de Datos}
\addcontentsline{toc}{section}{Apéndice C: Notas de Calidad de Datos}
\begin{itemize}
    \item Fecha de valoración: 20 de marzo de 2026
    \item Fuentes de tasas: cotizaciones mid de dealers
    \item Construcción de curvas: interpolación log-lineal entre nudos
    \item Ajuste de festivos: convención ACT/360 (no ajustada por festivos)
    \item Colateral: asumido respaldado en USD
\end{itemize}
